\documentclass[conference]{IEEEtran}
\usepackage{cite}
\usepackage{amsmath,amssymb,amsfonts}
\usepackage{graphicx}
\usepackage{textcomp}
\usepackage{xcolor}
\usepackage{booktabs}
\usepackage{multirow}
\usepackage{bm}
\usepackage{url}
\def\BibTeX{{\rm B\kern-.05em{\sc i\kern-.025em b}\kern-.08em
    T\kern-.1667em\lower.7ex\hbox{E}\kern-.125emX}}

\newcommand{\grad}{\nabla}
\newcommand{\divg}{\mathrm{div}}
\newcommand{\Real}{\mathbb{R}}
\newcommand{\loss}{\mathcal{L}}
\newcommand{\net}{u_{\theta}}
\DeclareMathOperator*{\argmin}{arg\,min}

\begin{document}

\title{Benchmarking PDE Priors in SIREN-Based PINNs\\
for Zero-Shot Single-Image Super-Resolution}

\author{%
\IEEEauthorblockN{1\textsuperscript{st} Alessandro Bottino}
\IEEEauthorblockA{\textit{University of Naples Federico II} \\
\textit{ISPP-CNR, Portici}\\
Naples, Italy}
\and
\IEEEauthorblockN{2\textsuperscript{nd} Antonio Marino}
\IEEEauthorblockA{\textit{University of Naples Federico II} \\
\textit{ISPP-CNR, Portici}\\
Naples, Italy}
\and
\IEEEauthorblockN{3\textsuperscript{rd} Antonio Vocino}
\IEEEauthorblockA{\textit{University of Naples Federico II}\\
Naples, Italy}
\and
\IEEEauthorblockN{4\textsuperscript{th} Andrea Solai}
\IEEEauthorblockA{\textit{University of Naples Federico II}\\
Naples, Italy}
\and
\IEEEauthorblockN{5\textsuperscript{th} Roberto Tessitore}
\IEEEauthorblockA{\textit{University of Naples Federico II}\\
Naples, Italy}
\and
\IEEEauthorblockN{6\textsuperscript{th} Andrea Salvatore}
\IEEEauthorblockA{\textit{University of Naples Federico II}\\
Naples, Italy}
\and
\IEEEauthorblockN{7\textsuperscript{th} Salvatore Cuomo}
\IEEEauthorblockA{\textit{University of Naples Federico II}\\
Naples, Italy}
\and
\IEEEauthorblockN{8\textsuperscript{th} Vincenzo Vocca}
\IEEEauthorblockA{\textit{University of Naples Federico II}\\
Naples, Italy}
}

\maketitle

\begin{abstract}
Single-image super-resolution can be written as an ill-posed inverse
problem in which many high-resolution images are consistent with the
same low-resolution observation. While supervised super-resolution
networks achieve strong visual quality, they require paired datasets
and encode the image prior implicitly. This work instead focuses on an
interpretable zero-shot setting and documents a benchmark built on
physics-informed neural networks (PINNs). The unknown high-resolution
image is represented by a continuous SIREN and constrained by a known
degradation operator composed of Gaussian blur and subsampling. Within
this unified framework we compare classical PDE-inspired priors,
namely Perona-Malik diffusion and vectorial anisotropic diffusion,
each coupled with either Total Variation or shock regularization. The
benchmark keeps the backbone, optimization schedule, warm start, and
synthetic degradation pipeline fixed so that only the regularization
strategy changes across trials. The final evaluation is designed on
the Set5 Baby and Butterfly images at $\times 2$ and $\times 4$, with
an additional small out-of-distribution check on up to two Darmstadt
images. This draft reports the full methodology, implementation
details, and result templates ready to be completed once the final
runs are available.
\end{abstract}

\begin{IEEEkeywords}
physics-informed neural networks, single-image super-resolution,
SIREN, anisotropic diffusion, total variation, benchmark
\end{IEEEkeywords}

\section{Introduction}
Single-image super-resolution (SISR) aims at recovering a
high-resolution (HR) image from a single low-resolution (LR)
observation. The problem is intrinsically ill-posed because the
forward operator discards high-frequency information, so infinitely
many HR images can explain the same LR input. Modern supervised
approaches such as SRCNN, EDSR, SwinIR, and diffusion-based models
obtain strong performance by learning an implicit prior from large
paired datasets \cite{dong2014srcnn,lim2017edsr,liang2021swinir,saharia2022sr3}.
However, they require offline training data and offer limited control
over the physical consistency of the reconstruction.

Classical PDE-based regularizers provide a complementary viewpoint.
Perona-Malik diffusion, anisotropic diffusion, Total Variation (TV),
and shock filters were originally introduced for denoising and
edge-preserving restoration \cite{peronamalik1990,weickert1998,rof1992,osherrudin1990}.
They are interpretable, modular, and well aligned with inverse
problems, but they are usually implemented on fixed grids and are less
commonly benchmarked inside continuous neural representations.

This paper documents a controlled benchmark where the HR image is
represented by a SIREN \cite{sitzmann2020siren} and regularized by
PDE-inspired losses inside a PINN framework \cite{raissi2019pinn}.
The emphasis is on benchmarking, not on proposing a new architecture.
The same backbone, degradation model, optimizer, and warm start are
kept fixed, while the regularizer is changed systematically. This
setup isolates the effect of scalar versus vectorial diffusion and of
TV versus shock sharpening in a zero-shot regime.

The main contributions of the proceeding are the following.
\begin{itemize}
\item A unified PINN formulation for zero-shot SISR with a continuous
      SIREN parameterization and a known blur-plus-subsampling forward
      operator.
\item A benchmark matrix that compares diffusion priors in two
      branches: diffusion plus TV and diffusion plus shock, under
      identical optimization settings.
\item A reproducible protocol on Set5 Baby and Butterfly at $\times 2$
      and $\times 4$, plus a small external check on two Darmstadt
      images used as additional HR sources before synthetic
      degradation.
\item A complete \LaTeX{} draft with sections, tables, and figure
      placeholders that can be filled directly once the final runs and
      training plots are produced.
\end{itemize}

The remainder of the paper is organized as follows. Section~II reviews
the closest literature. Section~III formulates the inverse problem.
Section~IV describes the PINN benchmark. Section~V details the
implementation choices. Section~VI defines the final protocol and the
result templates. Section~VII summarizes the planned analysis, and
Section~VIII concludes the paper.

\section{Related Work}
\subsection{Data-Driven Super-Resolution}
Deep learning has dominated SISR for the last decade. Early CNN models
such as SRCNN introduced direct end-to-end mapping from LR to HR
\cite{dong2014srcnn}; residual and transformer architectures such as
EDSR and SwinIR pushed the quantitative frontier further
\cite{lim2017edsr,liang2021swinir}; diffusion-based pipelines improved
perceptual realism \cite{saharia2022sr3}. These methods are highly
effective when a representative training distribution is available, but
their prior is learned offline and can be difficult to adapt when the
degradation operator is known and the data regime is extremely small.

\subsection{PDE-Based Image Restoration}
PDE-based models remain attractive because they encode restoration
assumptions explicitly. Perona-Malik diffusion suppresses smoothing
across strong gradients \cite{peronamalik1990}; anisotropic diffusion
extends this idea with directional control \cite{weickert1998}; TV
promotes sparse gradients and piecewise-smooth reconstructions
\cite{rof1992}; shock filtering sharpens edges through a nonlinear
transport mechanism \cite{osherrudin1990}. These priors were developed
mainly for denoising and deblurring, yet they are natural candidates
for the free subspace of an ill-posed SR inverse problem.

\subsection{PINNs and Implicit Representations}
PINNs enforce physical constraints through differentiable residuals
computed by automatic differentiation \cite{raissi2019pinn}. For image
restoration, they provide a convenient bridge between continuous
representations and PDE priors. SIREN networks are especially suited
to this setting because sinusoidal activations represent fine detail
and yield stable first- and second-order derivatives
\cite{sitzmann2020siren}. The repository used in this work also
contains a Fourier-feature MLP alternative \cite{tancik2020fourier},
but the benchmark described here fixes the backbone to SIREN in order
to isolate the role of the regularization terms.

\section{Problem Formulation}
\subsection{Forward Model}
Let $u \in [0,1]^{H \times W \times 3}$ be the unknown HR image and
$y \in [0,1]^{h \times w \times 3}$ the observed LR image. The
degradation model implemented in the codebase is
\begin{equation}
    y = P u + \eta, \qquad P = S B,
    \label{eq:forward}
\end{equation}
where $B$ is a separable Gaussian blur with standard deviation
$\sigma$, $S$ is stride-based subsampling with factor $s$, and $\eta$
is optional additive noise. In the current benchmark, the core SR
setting uses $\sigma = 1.0$, $s \in \{2,4\}$, and zero added noise so
that the comparison focuses on the regularizer rather than on noise
robustness.

\subsection{Inverse Problem}
The reconstruction is formulated as
\begin{equation}
    u^\star = \argmin_u \, \|P u - y\|_2^2 + \lambda R(u),
    \label{eq:energy}
\end{equation}
where $R(u)$ is a prior that biases the solution toward plausible HR
images. In the proposed PINN setting, $u$ is not optimized directly on
a discrete grid. Instead, it is represented by a neural function
$\net : [0,1]^2 \rightarrow \Real^3$ and the prior is enforced through
differentiable residual terms sampled at collocation points. This
allows the inverse problem to remain continuous while still using
classical PDE operators.

\section{Benchmark Method}
\subsection{Continuous Image Representation}
The benchmark fixes the HR image model to a SIREN
\cite{sitzmann2020siren}. Given normalized coordinates
$x \in [0,1]^2$, the network outputs RGB values
$\net(x) \in [0,1]^3$. The implementation uses five layers, hidden size
$256$, and first-layer frequency parameter $\omega_0 = 30$. This
representation is shared by all trials.

\subsection{Data Fidelity}
The main fidelity term in the repository is
\begin{equation}
    \loss_{\mathrm{data}}(\theta) =
    \|P \net - y\|_2^2,
    \label{eq:data}
\end{equation}
implemented by rendering the full HR grid, applying the same
blur-plus-subsampling operator used to synthesize the LR observation,
and computing the mean squared error. A second pointwise term exists in
the codebase, but it is not part of the core benchmark so that every
configuration is anchored by the same physically meaningful
$\loss_{\mathrm{data}}$.

\subsection{PDE and Sharpness Priors}
The benchmark compares two diffusion priors and two sharpening
regularizers.

\textbf{Perona-Malik diffusion.} For each channel,
\begin{equation}
    N_{\mathrm{PM}}[u] =
    \divg\!\left(g(\|\grad u\|^2)\grad u\right), \qquad
    g(s) = \frac{1}{1 + s/\kappa^2}.
    \label{eq:pm}
\end{equation}
This is the scalar edge-preserving baseline.

\textbf{Vectorial anisotropic diffusion.} RGB channels are coupled
through the structure tensor
\begin{equation}
    G(u) = \sum_{c=1}^{3} \grad u_c \grad u_c^\top.
    \label{eq:structure}
\end{equation}
Let $G = Q \operatorname{diag}(\lambda_1,\lambda_2) Q^\top$ and define
$D(u) = Q \operatorname{diag}(g_1,g_2) Q^\top$ with
$g_i = 1/(1+\lambda_i)$. The residual is
\begin{equation}
    N_{\mathrm{aniso}}[u]_c = \divg\!\left(D(u)\grad u_c\right).
    \label{eq:aniso}
\end{equation}
The current implementation evaluates $G(u)$ pointwise at the sampled
collocation coordinates and regularizes it before eigendecomposition.

\textbf{Total Variation.} The vectorial TV loss is
\begin{equation}
    \loss_{\mathrm{TV}}(\theta) =
    \mathbb{E}_x \left[
    \sqrt{\sum_{c=1}^{3} \|\grad u_c(x)\|^2 + \varepsilon}
    \right].
    \label{eq:tv}
\end{equation}
It promotes sparse gradients and therefore piecewise-smooth images
with sharper edges.

\textbf{Shock filter.} Let
$\eta = \grad u / (\|\grad u\| + \varepsilon)$ denote the local gradient
direction. The implemented differentiable shock residual is
\begin{equation}
    N_{\mathrm{shock}}[u] =
    \tanh\!\left(\frac{u_{\eta\eta}}{\varepsilon}\right)\|\grad u\|,
    \label{eq:shock}
\end{equation}
where $u_{\eta\eta}$ is the second directional derivative along
$\eta$. This term emphasizes edge sharpening but is potentially more
aggressive than TV.

\subsection{Boundary Conditions and Composite Loss}
Zero normal derivative boundary conditions are enforced on the four
borders of the normalized square domain. The total loss is
\begin{equation}
    \loss =
    \lambda_d \loss_{\mathrm{data}} +
    \alpha(t)\!\!\sum_{j \in \mathcal{P}}\!\!\lambda_j \loss_j +
    \lambda_{bc}\loss_{bc},
    \label{eq:total}
\end{equation}
where $\mathcal{P}$ contains the active prior terms and $\alpha(t)$ is
a linear curriculum that ramps PDE-type contributions from $0$ to $1$
after an initial warm-up. The benchmark will compare priors under this
same loss template.

\section{Implementation Details}
All experiments use the same settings taken from the current codebase.
HR images are resized to $256 \times 256$ before degradation. The
training loop runs for $2000$ epochs with Adam, learning rate
$10^{-4}$, $4096$ collocation points for prior terms, $2048$ sampled LR
data points, and gradient clipping at norm $1.0$. Before the main
optimization, the network is warm-started for $2000$ steps to reproduce
the bicubic upsampling of the LR image using Adam with learning rate
$10^{-3}$ and cosine annealing. This shared initialization is important
because it reduces poor local minima and makes cross-method comparisons
fairer.

Two numerical details are essential. First, derivatives are computed on
coordinates normalized to $[0,1]^2$, so second-order residuals scale
with the square of the image size. The implementation therefore divides
PDE residuals by $\max(H,W)^2$ to keep the loss magnitudes comparable
across resolutions. Second, the anisotropic tensor adds
$\varepsilon I$ before eigendecomposition to avoid numerical
instability in flat regions. The current reference weights are
$\lambda_{\mathrm{PM}} = \lambda_{\mathrm{aniso}} = 2 \times 10^{-2}$,
$\lambda_{\mathrm{TV}} = 2 \times 10^{-3}$,
$\lambda_{\mathrm{shock}} = 5 \times 10^{-3}$, and
$\lambda_{bc} = 10^{-3}$, with PDE warm-up starting after $50$ epochs
and ramping over the next $200$ epochs.

The repository also contains fourth-order terms and a learned CNN prior.
These components are deliberately excluded from the main benchmark
described in this paper, whose final comparison is restricted to the
SIREN backbone and to second-order diffusion priors combined with TV or
shock. Preliminary fourth-order sweeps are treated as exploratory only
and serve to motivate later follow-up experiments rather than the core
table of the proceeding.

\section{Experimental Protocol}
\subsection{Final Benchmark Matrix}
Table~\ref{tab:configs} defines the configurations planned for the
final benchmark. The matrix follows the latest project direction:
compare the same data term against two regularization branches,
diffusion plus TV and diffusion plus shock.

\begin{table}[t]
\caption{Planned benchmark configurations.}
\label{tab:configs}
\begin{center}
\small
\begin{tabular}{lll}
\toprule
ID & Active terms & Goal \\
\midrule
B0 & data & data-only PINN baseline \\
T0 & data + TV & sharpness-only baseline \\
T1 & data + PM + TV & scalar diffusion + TV \\
T2 & data + Aniso + TV & vectorial diffusion + TV \\
S0 & data + Shock & sharpness-only baseline \\
S1 & data + PM + Shock & scalar diffusion + shock \\
S2 & data + Aniso + Shock & vectorial diffusion + shock \\
\bottomrule
\end{tabular}
\end{center}
\end{table}

If the additional ZG-style candidate remains stable after the closing
exploratory runs, one extra row can be appended to each branch without
changing the rest of the protocol. In the codebase, these rows map to
\texttt{data\_lr}, \texttt{reg\_tv}, \texttt{pde\_perona\_malik},
\texttt{pde\_anisotropic}, and \texttt{pde\_shock}.

\subsection{Images, Scales, and Metrics}
The main evaluation is planned on the Set5 images \texttt{baby.png} and
\texttt{butterfly.png}. For each image, the same synthetic degradation
pipeline is applied and the methods are tested at scales
$s \in \{2,4\}$. Results will be reported in terms of PSNR and SSIM,
matching the metrics already implemented in the repository. The
comparison against bicubic interpolation is retained as the main
classical baseline, while the data-only PINN baseline B0 shows how much
is gained by the regularizer beyond the inverse constraint itself.

As a small out-of-distribution check, up to two additional images
originating from the Darmstadt benchmark will be used as external HR
sources before synthetic degradation. These extra runs are not intended
to define a new leaderboard; they only test whether the ranking of the
priors is stable outside the two Set5 reference images.

\subsection{Planned Outputs}
Besides the quantitative tables, the final paper will include:
\begin{itemize}
\item qualitative crops on Baby and Butterfly at $\times 2$ and
      $\times 4$;
\item training curves for data, PDE, and boundary terms, sampled every
      fixed number of epochs;
\item a short stability analysis comparing the convergence behavior of
      TV-based and shock-based branches.
\end{itemize}
The current trainer already produces snapshots every $25$ epochs, so
the loss-curve figure can follow the same cadence when the logging
script is completed.

\begin{table}[t]
\caption{Template for the final average results on Set5 Baby and Butterfly.}
\label{tab:main}
\begin{center}
\small
\begin{tabular}{lcc|cc}
\toprule
\multirow{2}{*}{Method} & \multicolumn{2}{c|}{$\times 2$} & \multicolumn{2}{c}{$\times 4$} \\
 & PSNR $\uparrow$ & SSIM $\uparrow$ & PSNR $\uparrow$ & SSIM $\uparrow$ \\
\midrule
Bicubic & -- & -- & -- & -- \\
B0 & -- & -- & -- & -- \\
T0 & -- & -- & -- & -- \\
T1 & -- & -- & -- & -- \\
T2 & -- & -- & -- & -- \\
S0 & -- & -- & -- & -- \\
S1 & -- & -- & -- & -- \\
S2 & -- & -- & -- & -- \\
\bottomrule
\end{tabular}
\end{center}
\end{table}

\begin{figure}[t]
\centering
\fbox{\parbox[c][1.25in][c]{0.47\columnwidth}{\centering
Qualitative crops\\
Baby and Butterfly\\
$\times 2$ and $\times 4$}}
\hfill
\fbox{\parbox[c][1.25in][c]{0.47\columnwidth}{\centering
Training curves\\
data, PDE, and BC losses\\
sampled every fixed epoch interval}}
\caption{Placeholders for the final qualitative and convergence
figures. The draft compiles without external assets and keeps the
layout ready for the final plots.}
\label{fig:placeholders}
\end{figure}

\section{Planned Analysis}
The benchmark is designed to answer a small set of focused questions.
First, does vectorial anisotropic diffusion improve over scalar
Perona-Malik diffusion when the RGB channels are coupled through a
shared structure tensor? This is particularly relevant near chromatic
edges, where channel-wise diffusion can produce color fringing.
Second, what is the trade-off between TV and shock regularization once
the same data term and the same diffusion prior are fixed? TV is
expected to be smoother and more stable, whereas shock-based variants
may recover crisper transitions at the price of higher sensitivity.

Third, how much do the gains depend on the difficulty of the inverse
problem? The $\times 2$ setting is relatively mild, while the
$\times 4$ setting leaves a larger null space to the prior. This makes
the pair $\{\times 2,\times 4\}$ informative even with only two core
images. Finally, the two Darmstadt-based runs will act as a small
sanity check on whether the qualitative ranking of the priors is tied
too strongly to Set5.

\section{Conclusion}
This draft consolidates the methodological part of a benchmark on
PDE-inspired priors for zero-shot single-image super-resolution with
SIREN-based PINNs. The paper now contains the full problem statement,
the exact benchmark matrix, the implementation details needed for
reproducibility, and the tables and figures required for the final
write-up. Once the last experimental runs are completed, the remaining
work is limited to filling the quantitative tables, inserting the
qualitative crops, and adding the convergence plots.

\begin{thebibliography}{00}
\bibitem{raissi2019pinn} M. Raissi, P. Perdikaris, and G. E. Karniadakis,
``Physics-informed neural networks: A deep learning framework for solving
forward and inverse problems involving nonlinear partial differential
equations,'' \textit{J. Comput. Phys.}, vol. 378, pp. 686--707, 2019.
\bibitem{sitzmann2020siren} V. Sitzmann, J. N. P. Martel, A. W. Bergman,
D. B. Lindell, and G. Wetzstein, ``Implicit neural representations with
periodic activation functions,'' in \textit{NeurIPS}, 2020.
\bibitem{tancik2020fourier} M. Tancik et al., ``Fourier features let
networks learn high frequency functions in low dimensional domains,''
in \textit{NeurIPS}, 2020.
\bibitem{peronamalik1990} P. Perona and J. Malik, ``Scale-space and edge
detection using anisotropic diffusion,'' \textit{IEEE Trans. Pattern
Anal. Mach. Intell.}, vol. 12, no. 7, pp. 629--639, 1990.
\bibitem{weickert1998} J. Weickert, \textit{Anisotropic Diffusion in Image
Processing}. Stuttgart: Teubner, 1998.
\bibitem{rof1992} L. I. Rudin, S. Osher, and E. Fatemi, ``Nonlinear total
variation based noise removal algorithms,'' \textit{Physica D}, vol. 60,
pp. 259--268, 1992.
\bibitem{osherrudin1990} S. Osher and L. I. Rudin, ``Feature-oriented image
enhancement using shock filters,'' \textit{SIAM J. Numer. Anal.}, vol. 27,
no. 4, pp. 919--940, 1990.
\bibitem{dong2014srcnn} C. Dong, C. C. Loy, K. He, and X. Tang,
``Image super-resolution using deep convolutional networks,''
\textit{IEEE Trans. Pattern Anal. Mach. Intell.}, vol. 38, no. 2,
pp. 295--307, 2015.
\bibitem{lim2017edsr} B. Lim, S. Son, H. Kim, S. Nah, and K. M. Lee,
``Enhanced deep residual networks for single image super-resolution,''
in \textit{CVPR Workshops}, 2017.
\bibitem{liang2021swinir} J. Liang et al., ``SwinIR: Image restoration
using Swin Transformer,'' in \textit{ICCV Workshops}, 2021.
\bibitem{saharia2022sr3} C. Saharia, J. Ho, W. Chan, T. Salimans,
D. J. Fleet, and M. Norouzi, ``Image super-resolution via iterative
refinement,'' \textit{IEEE Trans. Pattern Anal. Mach. Intell.}, 2022.
\bibitem{bevilacqua2012set5} M. Bevilacqua, A. Roumy, C. Guillemot, and
M.-L. Alberi-Morel, ``Low-complexity single-image super-resolution
based on nonnegative neighbor embedding,'' in \textit{BMVC}, 2012.
\bibitem{agustsson2017div2k} E. Agustsson and R. Timofte,
``NTIRE 2017 challenge on single image super-resolution: Dataset and
study,'' in \textit{CVPR Workshops}, 2017.
\end{thebibliography}

\end{document}
